\speciestitle{Methyl cyanide}{CH\cs3CN}{CH3CN}%
\begin{description}
\item[Useful range:] 46\,--\,1.0\,hPa
\item[Contact:] Michelle Santee, \email{Michelle.L.Santee@jpl.nasa.gov}
\end{description}
\label{sec:StandardCH3CN}%

\subsection*{Introduction}

The v2.2 standard CH\cs3CN data, which were derived from radiances
measured by the radiometer centered near 190\,GHz, were not recommended
for use in scientific studies.  In \thisv, the standard CH\cs3CN product
is taken from radiances measured by the radiometer centered near
640\,GHz.  In addition, the quality and reliability of the
\mbox{640-GHz} CH\cs3CN retrievals themselves have been improved in
\thisv, largely because of changes in the way that the continuum is being
handled for this radiometer in the v3 algorithms.  The MLS CH\cs3CN data
are now deemed scientifically useful over the range 46 to 1\,hPa, except
in the winter polar vortex regions, where they may exhibit large biases
below 10\,hPa.  In addition, the data at lower retrieval levels (i.e.,
higher pressures) may be used with caution in certain circumstances.  A
summary of the precision and resolution (vertical and horizontal) of the
\thisv CH\cs3CN measurements as a function of altitude is given in
Table~\ref{tab:CH3CNSummary}.  More details on the quality of the MLS
\thisv CH\cs3CN measurements are given below.

\subsection*{Resolution}
\onestandardavk{CH3CN}{CH\cs3CN}%

The resolution of the retrieved data can be described using ``averaging
kernels'' \citep[e.g.,][]{Rodgers00}; the two-dimensional nature of the
MLS data processing system means that the kernels describe both vertical
and horizontal resolution.  Smoothing, imposed on the retrieval system
in both the vertical and horizontal directions to enhance retrieval
stability and precision, degrades the inherent resolution of the
measurements.  Consequently, the vertical resolution of the \thisv
CH\cs3CN data, as determined from the full width at half maximum of the
rows of the averaging kernel matrix shown in Figure~\ref{fig:AvkCH3CN},
is \texttildelow5\,--\,6\,km in the lower stratosphere, degrading to
\texttildelow7\,--\,8\,km in the upper stratosphere.  Note that there is
overlap in the averaging kernels for the 100 and 147\,hPa retrieval
surfaces, indicating that the 147\,hPa retrieval does not provide as
much independent information as is given by retrievals at higher
altitudes.  Figure~\ref{fig:AvkCH3CN} also shows horizontal averaging
kernels, from which the along-track horizontal resolution is determined
to be \texttildelow400\,--\,700\,km over most of the vertical range.  The
cross-track resolution, set by the width of the field of view of the
\mbox{640-GHz} radiometer, is \texttildelow3\,km.  The along-track separation
between adjacent retrieved profiles is 1.5\degsym\ great circle angle
(\texttildelow165\,km), whereas the longitudinal separation of MLS
measurements, set by the Aura orbit, is 10\degsym--\,20\degsym\ over low
and middle latitudes, with much finer sampling in the polar regions.

\subsection*{Precision}

\begin{figure}
\centering\dontincludegraphics[width=14.0cm]{figures/CH3CN_ASCDSC_200804}
\caption{ (left) Ensemble mean profiles for ascending (red) and
  descending (blue) orbit matching pairs of MLS \thisv CH\cs3CN in the
  latitude range 50\degsym S\,--\,50\degsym N averaged over 30 days in April
  2008.  (center) Mean percent difference profiles between ascending and
  descending orbits (cyan), the standard deviation about the mean
  difference (orange), and the root sum square of the precisions
  calculated by the retrieval algorithm (magenta).  (right) Same, for
  absolute differences (pptv).  SD values are scaled by 1/$\sqrt{2}$;
  thus the observed SD represents the statistical repeatability of the
  MLS measurements, and the expected SD represents the theoretical
  \mbox{1-$\sigma$} precision for a single profile.  See
  \citet{LambertEtal07} for details.  }
\label{fig:ascdsc}
\end{figure}

The precision of the MLS CH\cs3CN data is estimated empirically by
computing the standard deviation of the differences between matched
measurement points at the intersections of the ascending (day) and
descending (night) sides of the orbit.  That the mean differences
between paired profiles are minimal (Figure~\ref{fig:ascdsc}) indicates
the absence of systematic ascending / descending biases.  Observed
scatter, representing the statistical repeatability of the measurements,
is 50\,--\,100\,pptv throughout the vertical domain.  This estimate reflects
the precision of a single profile; in most cases precision can be
improved by averaging, with the precision of an average of $N$ profiles
being 1/$\sqrt{N}$ times the precision of an individual profile (note
that this is not the case for averages of successive along-track
profiles, which are not completely independent because of horizontal
smearing).  The theoretical precision reported by the Level~2 data
processing system slightly exceeds the observationally-determined
precision throughout the vertical range, indicating that the smoothing
applied to stabilize the retrieval and improve the precision has a
nonnegligible influence.  Because the theoretical precisions take into
account occasional variations in instrument performance, the best
estimate of the precision of an individual data point is the value
quoted for that point in the L2GP files, but it should be borne in mind
that this approach slightly overestimates the actual measurement noise.

\subsection*{Range}

Although CH\cs3CN is retrieved (and reported in the L2GP files) over the
range 147 to 0.001\,hPa, on the basis of the drop off in precision and
resolution, the lack of independent information contributed by the
measurements, and the results of simulations using synthetic data as
input radiances to test the closure of the retrieval system, the data
are not deemed reliable at the extremes of the retrieval range.  Thus we
recommend that \thisv CH\cs3CN be used for scientific studies only at the
levels between 46 and 1\,hPa.  However, although the 147, 100, and
68\,hPa retrievals are not generally recommended, they may be
scientifically useful in some circumstances.  For example, the data
display unphysical sharp latitudinal gradients at $\pm$30\degsym\ at 100
and 68\,hPa, yet the large-scale longitudinal variations within the
tropics are probably robust.  Similarly, confined regions of significant
enhancement at 147\,hPa unaccompanied by comparably enhanced values at
100\,hPa may reflect real atmospheric features.  The \thisv CH\cs3CN data
at these levels (147\,--\,68\,hPa) should only be used in consultation with
the MLS science team.

\subsection*{Accuracy}

\begin{figure}
\centering
\dontincludegraphics[width=8.0cm]{figures/ch3cn_zonalmean_livesey2000}
\hspace{2em}
\parbox[b]{6cm}{%
\dontincludegraphics[width=5.8cm]{figures/Zm_CH3CN_v03-3x__2007d001-2007d031}\\
\dontincludegraphics[width=5.8cm]{figures/Zm_CH3CN_v03-3x__2007d182-2007d212}}
\caption{(left plot) Top row shows UARS MLS mean CH\cs3CN fields for
  June / July 1993 (left) and December 1992 / January 1993 (right).  The
  other rows show results from various chemistry transport model runs.
  See \citet{LiveseyEtal01} for details. (right plot) \thisv Aura MLS
  CH\cs3CN monthly zonal means for January (top) and July (bottom)
  2007. }
\label{fig:ch3cnZM}
\end{figure}

The impact of various sources of systematic uncertainty has not yet been
quantified for CH\cs3CN as it has for most other MLS products.  This
work is planned as part of a dedicated validation exercise for the \thisv
CH\cs3CN data.  However, preliminary comparisons with results from a
two-dimensional chemistry transport model and CH\cs3CN retrievals from
the MLS instrument on the Upper Atmosphere Research Satellite (UARS)
\citep{LiveseyEtal01} indicate that the \thisv Aura MLS CH\cs3CN mixing
ratios are biased substantially high in the lower stratosphere
(147\,--\,68\,hPa, see Figure~\ref{fig:ch3cnZM}).  Furthermore, the
zonal-mean morphology of the Aura MLS CH\cs3CN at the lowest levels does
not agree well with that either observed by UARS MLS or predicted by the
model.

\subsection*{Review of comparisons with other datasets}

Detailed comparisons with correlative data sets have not yet been
undertaken.  This work is planned as part of a dedicated validation
exercise for the \thisv CH\cs3CN data.

\subsection*{Data screening}
\begin{description}
\item[Pressure range:] \screeningrule{46\,--\,1.0\,hPa}
    \standardrangestatement\ The CH\cs3CN data at 147\,--\,68\,hPa may be
    useful under certain circumstances but should not be analyzed in
    scientific studies without significant discussion with the MLS
    science team.

\standardprecisionrule

\evenstatusrule

\item[Clouds:] \cloudsokrule

  Nonzero but even values of \texttt{Status} indicate that the profile
  has been marked as questionable, usually because the measurements may
  have been affected by the presence of thick clouds.  Globally fewer
  than \texttildelow1\,--\,2\% of CH\cs3CN profiles are typically identified in
  this manner (though this value rises to \texttildelow3\,--\,5\% in the tropics on
  a typical day), and clouds generally have little influence on the
  stratospheric CH\cs3CN data.  Thus profiles with even values of
  \texttt{Status} may be used without restriction.

\item[Quality:] \qualityrule{1.4}

  This threshold for \texttt{Quality} (see
  section~\ref{sec:StatusQuality}) typically excludes less than 1\% of
  CH\cs3CN profiles on a daily basis; note that it potentially discards
  some ``good'' data points while not necessarily identifying all
  ``bad'' ones.
\item[Convergence:] \convergencerule{1.05}

  On a typical day this threshold for \texttt{Convergence} (see
  section~\ref{sec:StatusQuality}) discards very few (0.3\% or less) of
  the CH\cs3CN profiles, many (but not all) of which are filtered out by
  the other quality control measures.
\end{description}

\subsection*{Artifacts}
\begin{itemize}
\item The retrievals at 100 and 68\,hPa are characterized by unphysical
  sharp latitudinal gradients at $\pm$30\degsym.
\item Substantial biases may be present in the mixing ratios in the
  winter polar vortex regions for retrieval levels in the range
  100\,--\,15\,hPa.
\end{itemize}

\subsection*{Desired improvements for future data version(s)}
\begin{itemize}
\item Improve the CH\cs3CN retrievals at 147\,--\,68\,hPa.
\end{itemize}

\begin{table*}[bp]
\caption{Summary of Aura MLS \thisv CH\cs3CN Characteristics}\vspace{0.5ex}
\label{tab:CH3CNSummary}
\begin{minipage}{\linewidth}
\renewcommand{\thefootnote}{\thempfootnote}
\centering\begin{tabular}{ccccl}
\toprule
\parbox[m]{18mm}{\bfseries\centering{Pressure\\ / hPa}} &
\parbox[m]{24mm}{\bfseries\centering{Resolution\\ V $\times$ H%
\footnote{Vertical and Horizontal resolution in along-track direction.} \\ / km}} &
\parbox[m]{18mm}{\bfseries\centering{Precision%
    \footnote{Precision on individual profiles.} \\/ pptv}} &
\parbox[m]{20mm}{\bfseries\centering{Accuracy\\ / pptv}} &
\parbox[m]{40mm}{\bfseries\centering{Known Artifacts\\ or Other Comments}} \\
%
\midrule
%
0.68\,--\,0.001 & ---   & ---    & ---     & Unsuitable for scientific
  use\\
%
1.0      & 6 $\times$ 800 & $\pm$100 & TBD & \\
%
46\,--\,1.5  & 5\,--\,8 $\times$ 400\,--\,700 & $\pm$50 & TBD & \\
%
100\,--\,68  & 5\,--\,6 $\times$ 600\,--\,700 & $\pm$50 & TBD & Consult with
MLS science team \\
%
147      & 4 $\times$ 800        & $\pm$100 & TBD & Consult with
MLS science team \\
%
1000\,--\,215  & ---    & ---     & ---      & Not retrieved\\
%
\bottomrule
\end{tabular}
\end{minipage}
\end{table*}

%%% Local Variables: 
%%% mode: latex
%%% TeX-master: "v4-x-report"
%%% End: 
